% !TEX root = 00_main.tex
% TLDR: The judge is the research object; VeriBench supplies prior context and artifacts.
\section{Introduction}
% CS197 move 1: problem motivation -- incomplete statements can accompany valid proofs.
A formalization can prove every theorem it states and still omit a required property of the source program.
Assessing theorem coverage therefore requires comparing the candidate statements with the intended obligations, not merely checking proof completion.
We call this property \emph{formal specification coverage} (FSC).
The judge is the object of this paper: given a source program, candidate formalization and reference, how well does it predict human ratings?
\textsc{VeriBench} provides the application and historical artifacts for this question~\citep{veribench2026}.
Its broader code-verification benchmark is related work rather than this paper's contribution.

% CS197 move 2: existing practice -- agreement evidence is useful but task-specific.
Large language model (LLM) judges offer a practical way to compare outputs with instructions and references.
\citet{zheng2023judging} study agreement with human preferences and document position, verbosity and self-enhancement biases.
Such agreement depends on the task, rating procedure and judge configuration.
In our historical FSC data, strong rank correlation coexists with a positive mean rating bias (Section~\ref{sec:legacy}).
A ranking statistic alone cannot identify the human rating distribution for an individual artifact.
This observation motivates a richer target; it does not establish that correlation training performs worse than loss training.

% CS197 move 3: contribution -- condition on each artifact and retain human disagreement.
We model human ratings as a conditional distribution over the six categories zero through five.
The judge receives the original Python source, candidate Lean artifact and reference shown to the raters.
It predicts a probability for each rating rather than only a single scalar rating.
Keeping individual ratings preserves information about disagreement that a consensus label would discard.
Proper scoring rules provide losses whose expected optimum is the target distribution~\citep{gneiting2007strictly}.
Finite data, restricted judge classes and optimization can still prevent a fitted judge from attaining that optimum.

% CS197 move 4: instantiation -- separate training objective from validation selection.
We specify a comparison between training with Pearson correlation and training with a proper probability loss.
Both conditions use the same judge class, training data, computational budget and validation criterion.
Normalized ranked probability score (RPS) is the primary training loss for the proper-loss condition and the common validation criterion.
Cross-entropy and Brier loss define additional, separately reported training conditions.
The fixed judges share a held-out evaluation panel with RPS primary and cross-entropy and total variation distance secondary.
Pearson correlation measures linear association of output ratings; optimizing it does not require a linear judge.

% CS197 move 5: evidence -- historical samples, current collection and sizing are separate records.
The study has three connected parts: historical judge evidence, new human-rating collection, and estimation of the required collection size.
% Provenance: experiments/07_exp1_metrics/results.md, section 1; original source be287d9.
The historical pool contains 25 tasks with three candidate samples each, yielding 75 artifacts.
Five rater columns contain 345 observed ratings out of 375 possible; the primary three-expert subset contains 195 out of 225.
The three candidates are samples, not three solver configurations.
The new collection prepares versioned candidate packets and records each person's rating under a common rubric.
Its main-round judge comparison is pending, and the final larger-study allocation remains open.
The historical results motivate these choices without supplying independent validation for the new judge.

% CS197 move 6: implications and contributions -- empirical question remains open.
The intended outcome is a judge whose predictive error and transfer limits can be measured against a specified human-rating target.
Human agreement remains distinct from semantic truth and from a checked implication proof.
We make the following contributions:
\begin{enumerate}
\item We define an artifact-conditional ordinal prediction task and a matched comparison of correlation training with proper-loss training.
\item We consolidate historical judge evidence with its rating counts, input provenance, uncertainty and prior-exposure limitations.
\item We specify how the ongoing collection and sample-size study will support a prospective comparison without treating planned ratings as completed evidence.
\end{enumerate}
